% concatenated from 2026.02.06.weighted-soliton-csck.tex
\documentclass[12pt]{amsart}

\usepackage[margin=30truemm]{geometry}
\usepackage{amssymb} 
\usepackage{braket}
\usepackage{bm}
\usepackage{comment}
\usepackage{mathtools}

%\usepackage{showkeys}%後で取る
\setcounter{tocdepth}{1}%目次レベル
\numberwithin{equation}{section}%章ごとの式番号
\mathtoolsset{showonlyrefs=true}%引用式だけ番号付け

\renewcommand{\baselinestretch}{1.2}

\theoremstyle{plain} 
\newtheorem{theorem}{Theorem}[section]
\newtheorem{lemma}[theorem]{Lemma}
\newtheorem{corollary}[theorem]{Corollary}
\newtheorem{proposition}[theorem]{Proposition}
\newtheorem{claim}[theorem]{Claim}
\newtheorem{remark}[theorem]{Remark}

\theoremstyle{definition}
\newtheorem{definition}[theorem]{Definition}
\newtheorem{example}[theorem]{Example}
\newtheorem{notation}[theorem]{Notation}
\newtheorem{assertion}[theorem]{Assertion}
\newtheorem*{remark0}{Remark}
\newtheorem*{ackn}{\bf{Acknowledgment}}
\newtheorem*{orga}{\bf{Organization}}

\renewcommand{\proofname}{Proof.} 

%%%%%%%%%    Original Notation    %%%%%%%%%%%%%
\newcommand{\Vol}[0]{\operatorname{Vol}}
\newcommand{\Pic}[0]{\operatorname{Pic}}
\newcommand{\Sym}[0]{\operatorname{Sym}}
\newcommand{\End}[0]{\operatorname{End}}
\newcommand{\rank}[0]{\operatorname{rank}}
\newcommand{\codim}[0]{\operatorname{codim}}
\newcommand{\sign}[0]{\operatorname{sign}}
\newcommand{\Coker}[0]{\operatorname{Coker}}
\newcommand{\Ker}[0]{\operatorname{Ker}}
\newcommand{\Image}[0]{\operatorname{Im}}
\newcommand{\Supp}[0]{\operatorname{Supp}}

\newcommand{\deldelb}{\sqrt{-1}\partial_{B} \overline{\partial_{B}}}
\newcommand{\deldel}{\sqrt{-1}\partial \overline{\partial}}
\newcommand{\del}{\sqrt{-1}\partial}
\newcommand{\dbar}{\overline{\partial}}
\newcommand{\e}{\varepsilon}
\newcommand{\ome}{\widetilde{\omega}}
\newcommand{\I}[1]{\mathcal{J}(#1)}
\newcommand{\mi}{{\rm{min}}}
\newcommand{\sing}{{\rm{sing}}}
\newcommand{\lla}[0]{{\langle\!\hspace{0.02cm} \!\langle}}
\newcommand{\rra}[0]{{\rangle\!\hspace{0.02cm}\!\rangle}}
\newcommand{\poly}{\bigtriangleup}
\newcommand{\ke}{\omega_{\rm KE}}
\newcommand{\tr}{\mathrm{tr}}
\newcommand{\Ric}{\mathrm{Ric}}
\newcommand{\MA}{\mathrm{MA}}
\newcommand{\lext}{l^{\mathrm{ext}}}

\renewcommand{\a}{\alpha}
\renewcommand{\b}{\beta}
\newcommand{\g}{\gamma}
\renewcommand{\d}{\delta}
%\newcommand{\e}{\varepsilon}
\newcommand{\s}{\sigma}
\renewcommand{\t}{\theta}
%\renewcommand{\l}{\lambda}
\renewcommand{\r}{\rho}
%\renewcommand{\o}{\omega}
\renewcommand{\phi}{\varphi}

\newcommand{\tA}{\tilde{A}}

\renewcommand{\ae}{{\rm a.e.\ }} 
\newcommand{\cf}{{\rm cf.\ }} 
\newcommand{\eg}{{\rm e.g.\ }} 
\newcommand{\ie}{{\rm i.e.\ }} 

\renewcommand{\i}{\sqrt{-1}}
\renewcommand{\leq}{\leqslant}
\renewcommand{\geq}{\geqslant}
\newcommand{\abs}[1]{\left\lvert#1\right\rvert}
\newcommand{\norm}[1]{\left\|#1\right\|} 
\newcommand{\ip}[1]{\left\langle#1\right\rangle} 
\newcommand{\iip}[1]{\left\lbrace#1\right\rbrace} 
\newcommand{\1}{\mathbf{1}}
%%%%%%%%%%%%%%%%%%%%%%%%%%%%%%%%%%%%%%%%

\begin{document}

\title[Soliton-type metrics associated with weighted CSCK metrics]
{Soliton-type metrics associated with weighted CSCK metrics
on Fano manifolds}
\author[S. Nakamura]
{Satoshi Nakamura} 

\subjclass[2010]{ %2010 MSC numbers
Primary 53C25; Secondary 53C55, 58E11.
}
\keywords{ %key words and phrases
$(v,w)$-CSCK metrics, $g$-solitons, Fano manifolds, Sasaki manifolds
}
\thanks{ %Thanks
}
\address{
}
\email{}

\address{
S. Nakamura:
Department of Mathematics, Institute of Science Tokyo,
2-12-1, Ookayama, Meguro-ku, Tokyo, 152-8551, Japan
}
\email{s.nakamura@math.titech.ac.jp}

\date{\today}

%%%%%%%%%%%%%%%%%%%%%%%%%%%%%%%%%%%%%%%%%

\maketitle

\begin{abstract}
We study weighted constant scalar curvature K\"ahler metrics,
introduced by Lahdili as $(v,w)$-CSCK metrics,
on Fano manifolds and their relationship with soliton-type metrics.
In this paper, we introduce a weight function $g(v,w)$
associated with a pair of weight functions $(v,w)$.
Assuming that $v$ and $g(v,w)$ are positive and log-concave
on the moment polytope,
we prove that the existence of a $(v,w)$-CSCK
metric in the first Chern class is equivalent to the existence of a
$g(v,w)$-soliton.

We also explain that a $g(v,w)$-soliton arises naturally from
Sasaki geometry.
More precisely, let $(v,w)$ be the weight functions defining a weighted CSCK metric
in $2\pi c_1(X)$ which gives rise to a $\hat{\xi}$-transverse extremal
metric on an $S^1$-bundle $N$ in the canonical bundle of a Fano
manifold $X$, where $\hat{\xi}$ is a possibly irregular Reeb field on $N$.
We prove that the associated $g(v,w)$-soliton on $X$ gives rise to a
$\hat{\xi}$-transverse Mabuchi soliton on $N$.
\end{abstract}


%\newpage
\tableofcontents


%%%%%%%%%%%%%%%%%%%%%%%%%%%%%%%%%%%%%
\section{Introduction}
In his paper \cite{Lah19}, Lahdili
introduced the notion of a weighted
constant scalar curvature K\"ahler metric (a weighted  CSCK metric for short)
formulated as follows.
Let $X$ be an $n$-dimensional compact K\"ahler manifold with a
K\"ahler class $\Omega$, and 
let $T$ be a maximal compact torus in
the reduced automorphism group $\mathrm{Aut}_{r}(X)$.
Fix two smooth functions $v>0$ and $w$ defined on
the moment map image $P_{X}$.
A $T$-invariant K\"ahler metric $\omega\in\Omega$ is called
 a $(v,w)$-constant scalar curvature K\"ahler metric
(a $(v,w)$-CSCK metric for short) if $S_{v}(\omega)=w(m_{\omega})$,
where $S_{v}(\omega)$ is the $v$-weighted scalar curvature of $\omega$ and 
$m_{\omega}$ is the moment map for the $T$-action
with respect to $\omega$.
When $v\equiv 1$, the weighted scalar curvature equals the ordinary scalar curvature.
The $(v,w)$-CSCK metrics provide a unifying framework that encompasses several
well-known classes of canonical metrics,
including
constant scalar curvature K\"ahler metrics, extremal K\"ahler metrics, 
extremal Sasaki metrics, and so on,
as reviewed in Section \ref{weighted section}.
In recent years, $(v,w)$-CSCK metrics itself have attracted considerable attention. 
The existence of $(v,w)$-CSCK metrics with log-concave weight $v$
was shown by Han-Liu \cite{HanLiu25} and Di~Nezza-Jubert-Lahdili \cite{DJL24, DJL25}
to be equivalent to the coercivity of the associated Mabuchi-type energy functional. 
Moreover a Yau-Tian-Donaldson-type correspondence was established
very recently by Boucksom-Jonsson \cite{BJ25},
building on earlier work of Li \cite{CL22}
(see also \cite{HanLiu25b, Ha25}),
characterizing the existence of $(v,w)$-CSCK metrics with log-concave weight $v$
in terms of an appropriate notion of K-stability.

In the case where $X$ is a Fano manifold and $\Omega=2\pi c_{1}(X)$,
soliton-type metrics called $g$-solitons,
where $g$ is a smooth positive function defined on $P_{X}$,
play a central role within the class of weighted CSCK metrics.
A $T$-invariant metric $\omega\in 2\pi c_{1}(X)$ is called a $g$-soliton
if $\Ric_{g}(\omega)=\omega$, where $\Ric_{g}(\omega)$ is the $g$-weighted
Ricci curvature form of $\omega$.
When $g\equiv 1$, the wighted Ricci form equals the ordinary Ricci form.
A $g$-soliton equals a $(g,w)$-CSCK metric with
$w(x)=n+\langle (\log g)^{\prime}(x),x\rangle$ for $x\in P_{X}$ \cite{AJL23}.
The notion of $g$-solitons provides a unified framework encompassing
K\"ahler-Einstein metrics, K\"ahler-Ricci solitons, Mabuchi solitons, Sasaki-Einstein metrics,
and so on, as reviewed in Section \ref{weighted section}.
Originating in Mabuchi's work \cite{Ma03} and further developed 
by Berman-Witt Nysr\"om \cite{BW14} and Han-Li \cite{HaLi20},
the existence of $g$-solitons with the log-concave weight $g$
is by now well understood to be characterized by
the coercivity of the associated Ding-type functional
as well as by an appropriate Ding-type stability condition.

Recently, there has been growing progress in the study of relationships
between the existence of two, in general different, canonical K\"ahler metrics
on a Fano manifold.
For instance, it was proved by Apostolov-Lahdili-Nitta \cite{ALN25} and
by Hisamoto and the author \cite{HN24} (see also \cite[Section 9.6]{Ma21}) that
the existence of a $T$-invariant Calabi's extremal metric \cite{Ca82} in $2\pi c_{1}(X)$
whose scalar curvature is strictly less than $n+1$
is equivalent to that of a $T$-invariant Mabuchi's soliton \cite{Ma01} in $2\pi c_{1}(X)$.
Here an extremal metric means a $(1,n+l_{\mathrm{ext}})$-CSCK metric
for an affine function $l_{\mathrm{ext}}$ on $P_{X}$
corresponds to the extremal vector field \cite{FM95} on $X$,
and a Mabuchi soliton means a $(1-l_{\mathrm{ext}})$-soliton.
More generally we already know the following.

\begin{theorem}{\rm (\cite{ALN25, HN24}, see also \cite{NN25})}\label{meta theorem}
Let $X$ be an $n$-dimensional Fano manifold.
Assume the positive weight function $g$ is log-concave on $P_{X}$.
Then the following conditions are equivalent.
\begin{enumerate}
\item There exists a $T$-invariant $(1, 1+n-g)$-CSCK metric
(i.e. $g$-extremal metric in the sense of \cite{HN24, NN25})
in $2\pi c_{1}(X)$.
\item There exists a $T$-invariant $g$-soliton in $2\pi c_{1}(X)$.
\end{enumerate}
\end{theorem}

The main result of this article is an extension of Theorem \ref{meta theorem},
formulated within the framework of $(v,w)$-CSCK metrics
on Fano manifolds, which we state as follows.

\begin{theorem}\label{main theorem}
Let $X$ be an $n$-dimensional Fano manifold.
Assume the positive weight function $v$ is log-concave on $P_{X}$.
Assume also the associated weight function
\begin{equation}\label{g(v,w)}
g(v,w)(x):=v(x)(1+n+\langle(\log v)^{\prime}(x),x\rangle -w(x))
\end{equation}
is positive and log-concave on $P_{X}$.
Then the following conditions
are equivalent.
\begin{enumerate}
\item There exists a $T$-invariant $(v,w)$-CSCK metric in $2\pi c_{1}(X)$.
\item There exists a $T$-invariant $g(v,w)$-soliton in $2\pi c_{1}(X)$.
\end{enumerate}
\end{theorem}

In particular, Theorem \ref{main theorem} shows that
the existence of a $(v,w)$-CSCK metric on a Fano manifold is governed by
the existence of a soliton-type metric,
namely the $g(v,w)$-soliton.

In view of Theorem \ref{meta theorem}, one obtains the following.

\begin{corollary}\label{sub theorem}
Under the same assumptions as in Theorem \ref{main theorem},
the following conditions are equivalent.
\begin{enumerate}
\item There exists a $T$-invariant $(v,w)$-CSCK metric in $2\pi c_{1}(X)$.
\item There exists a $T$-invariant $(1, 1+n-g(v,w))$-CSCK metric in $2\pi c_{1}(X)$.
\end{enumerate}
\end{corollary}

\begin{remark}
The author has been informed by V.~Apostolov and Y.~Nitta that
the same results as Theorem \ref{main theorem} and Corollary \ref{sub theorem} 
were obtained independently in a joint work with Apostolov-Lahdili-Nitta \cite{ALN25b}.
\end{remark}

One piece of evidence for the equivalence between the existence of $(v,w)$-CSCK metrics and
that of $g(v,w)$-solitons is that the Futaki-type invariant
$M_{v,w}^{\prime}$ of $X$, which provides an obstruction to the existence of 
$(v,w)$-CSCK metrics, coincides with another Futaki-type invariant 
$D_{g(v,w)}^{\prime}$ of $X$, which obstructs the existence of $g(v,w)$-solitons.
See Proposition \ref{FutMD}.
With a suitable formulation of the weight function $g(v,w)$,
the arguments developed by Hisamoto and the author \cite{HN24} can be extended.
A crucial step in our proof of Theorem \ref{main theorem}
is to establish Theorem \ref{bound of M} which shows
the uniform boundedness of the Mabuchi-type functional $M_{v,w}$
associated with $(v,w)$-CSCK metrics
along a continuity path used to construct a $g(v,w)$-soliton.

In Theorem \ref{main theorem},
we assume the positivity and log-concavity of the weight functions $v$ and $g(v,w)$.
The positivity assumption is essential in order for the weighted Monge-Amp\`ere measure
to be positive.
On the other hand, the log-concavity assumption is used in the regularity argument for weak solutions
\cite{DJL24, DJL25, HaLi20, HanLiu25}.
At present, it is not clear to the author whether this log-concavity assumption can be relaxed.

One of the main significances of Theorem \ref{main theorem} is that the existence of 
$(v,w)$-CSCK metrics on Fano manifolds can be captured by the existence of 
$g(v,w)$-solitons, that is, by the solvability of a complex Monge-Amp\`ere equation,
and hence can be characterized by a Ding-type stability condition,
which are deeply connected with the theory of minimal model program
in birational geometry \cite{HaLi20}.
The author expects that Theorem \ref{main theorem}
may be applied to the moduli theory of Fano manifolds admitting $(v,w)$-CSCK metrics.

An important and interesting application of $(v,w)$-CSCK metrics and $g$-solitons
is that they allow one to describe canonical Sasaki metrics on a Sasaki manifold
with a possibly irregular Reeb field
via (quasi-) regular quotients \cite{ACL21, AJL23}.
Let $X$ be a Fano manifold and $\omega_{0}\in 2\pi c_{1}(X)$ a K\"ahler metric
invariant under the action of a maximal compact torus $T$.
Let $N$ be the Sasaki manifold with the regular Reeb filed $\hat{\chi}$
associated with $(X,\omega_{0})$.
Namely $\pi:N_{\omega_{0}}\to X$ is the $S^{1}$-bundle contained in the canonical bundle $K_{X}$
with respect to the Hermitian metric on $K_{X}$ whose curvature is $-\omega_{0}$,
and $\hat{\chi}$ is the generator of the $S^{1}$-action.
Then one obtains a strongly pseudo-convex CR manifold $(N_{\omega_{0}}, D_{0},J_{0})$,
where $D_{0}$ is the horizontal distribution and $J_{0}$ is induced complex structure on $D_{0}$.
Fix the new torus $\hat{T}:=T\times S^{1}$ acting on $N_{\omega_{0}}$,
where $T$ acts on $K_{X}$ as a lift and $S^{1}$ is generated by $\hat{\chi}$.
Let $\xi\in\mathrm{Lie}(T)\subset\mathrm{Lie}(\hat{T})$
and consider $\hat{\xi}:=\xi+\hat{\chi}\in\mathrm{Lie}(\hat{T})$.
Now we assume that $\hat{\xi}$ defines a Reeb field for
$(N_{\omega_{0}}, D_{0},J_{0})$.
One has the corresponding contact $1$-form $\eta_{0}^{\hat{\xi}}$ on $N_{\omega_{0}}$
characterized by the conditions
$\eta_{0}^{\hat{\xi}}(\hat{\xi})=1$ and $\Ker\eta_{0}^{\hat{\xi}}=D_{0}$.
In these setting, as observed in Section \ref{Sasaki section},
the defining equation of a $\hat{\xi}$-transverse
extremal metric \cite{BGS08} in the basic cohomology class
$[d\eta_{0}^{\hat{\xi}}]_{B,\hat{\xi}}$,
which is in fact equal to the basic first Chern class
$2\pi c_{1}^{B,\hat{\xi}}(N_{\omega_{0}})$,
can be rewritten as the defining equation of a weighted CSCK metric
on $X$ in $2\pi c_{1}(X)$ for an appropriate weights $(v_{\xi},w_{\xi})$,
and conversely.
Similarly, the defining equation of a $\hat{\xi}$-transverse Mabuchi soliton
in $[d\eta_{0}^{\hat{\xi}}]_{B,\hat{\xi}}$
can also be rewritten as that of an appropriate soliton-type metric on $X$,
and conversely.

As shown in the following result,
an example of a $g(v,w)$-soliton
appears naturally in Sasaki geometry of $N_{\omega_{0}}$
for the possibly irregular Reeb field $\hat{\xi}$.
Moreover as observed in Section \ref{Sasaki section} 
this $g(v,w)$-soliton is associated with
weights satisfying $v(x)\not\equiv 1$ and
$w(x)\neq n+\langle (\log v)^{\prime}(x),x\rangle$
in general.

\begin{theorem}\label{Sasaki theorem}
Let $(v_{\xi},w_{\xi})$ be the weights of a weighted CSCK metric
in $2\pi c_{1}(X)$ giving rise to a $\hat{\xi}$-transverse 
extremal metric in $[d\eta_{0}^{\hat{\xi}}]_{B,\hat{\xi}}$.
Then the $g(v_{\xi}, w_{\xi})$-soliton in $2\pi c_{1}(X)$
gives rise to a $\hat{\xi}$-transverse Mabuchi soliton
in $[d\eta_{0}^{\hat{\xi}}]_{B,\hat{\xi}}$.
\end{theorem}

We shall give a proof from a general viewpoint of
a $\hat{\xi}$-transversal $g$-extremal metric and 
a $\hat{\xi}$-transversal $g$-soliton.
See Theorem \ref{Sasaki g(v,w)}.

As an extension of the K\"ahler case, 
it is a natural problem to ask whether
the existence of a transverse extremal metric
whose transverse scalar curvature is strictly less than $n+1$
is equivalent to that of a transverse Mabuchi soliton. 
Unfortunately Theorem \ref{main theorem} and \ref{Sasaki theorem}
do not currently provide an answer to this question,
since the weight functions $v_{\xi}$ and $g(v_{\xi},w_{\xi})$
are not log-concave in general.
However recall that the log-concavity assumption on the weights
in Theorem \ref{main theorem} is a technical condition
used to obtain regularity of weak solutions
\cite{DJL24, DJL25, HaLi20, HanLiu25}.
Therefore Theorem \ref{main theorem} and \ref{Sasaki theorem} provide strong 
evidence supporting the validity of the above,
as yet unproven, equivalence.

\begin{orga}
In Section~2, we recall basic notions on weighted CSCK metrics and
soliton-type metrics, including weighted scalar curvature and
weighted Ricci curvature.
Typical situations in which these metrics appear are also reviewed.
Section~3 introduces $g(v,w)$-solitons and explains their relationship
with $(v,w)$-CSCK metrics from the viewpoint of Futaki-type invariants
and weighted energy functionals.
In Section~4, we prove Theorem \ref{main theorem} by combining
coercivity properties of weighted energy functionals with a continuity
method for $g(v,w)$-solitons.
Finally, in Section~5, 
we prove Theorem \ref{Sasaki theorem}
(and Theorem \ref{Sasaki g(v,w)} for more general viewpoint).
For this, we discuss
how $(v,w)$-CSCK metrics and $g(v,w)$-solitons arise from
Sasaki-extremal metrics and Sasaki-Mabuchi solitons via regular quotients.
\end{orga}

\begin{ackn}
The author would like to thank Vestislav Apostolov and Yasufumi Nitta for sharing
their research article and for helpful discussions and comments.
He would like to thank Tomoyuki Hisamoto for many stimulating discussions.
He also would like to thank Eiji Inoue and Abdellah Lahdili for helpful comments.
He is supported by JSPS Grant-in-Aid for Early-Career Scientists
No.~24K16917.
\end{ackn}
%%%%%%%%%%%%%%%%%%%%%%%%%%%%%%%%%%%%%
\section{Preliminaries on $(v,w)$-CSCK metrics and $g$-solitons}
\label{weighted section}
In this section we fix some notation to give precise definitions of
weighted CSCK metrics and soliton-type metrics.
Let $X$ be an $n$-dimensional compact K\"ahler manifold with a
K\"ahler class $\Omega$.
Let $T$ be a maximal compact torus in
the reduced automorphism group $\mathrm{Aut}_{r}(X)$,
and $\frak{t}$ be the Lie algebra of $T$.
Fix a $T$-invariant metric $\omega_{0}\in\Omega$ to define
the space of $T$-invariant K\"ahler potentials $\mathcal{H}^{T}$
with respect to $\omega_{0}$.
For any $\phi\in\mathcal{H}^{T}$,
the $T$-invariant metric $\omega_{\phi}:=\omega_{0}+\deldel\phi$
gives rise to the moment map
\begin{equation}
m_{\phi}:X\to\mathfrak{t}^{*}
\end{equation}
defined by the equation
$i_{\xi}\omega_{\phi}=-d\langle m_{\phi},\xi \rangle$ and
the normalization $\int_{X}\langle m_{\phi},\xi \rangle\omega_{\phi}^{n}=0$
for any $\xi\in\frak{t}$.
It is well known that the image $P_{X}:=m_{\phi}(X)$ is
a convex polytope
independent of the choice of $\phi\in\mathcal{H}^{T}$.
Fix a strictly positive and smooth function $v$ on $P_{X}$.
One can assume that
the weighted Monge-Amp\`ere measure
\begin{equation}
\MA_{v}(\phi):=\frac{v(m_{\phi})\omega_{\phi}^{n}}{\int_{X}\omega_{0}^{n}}
\end{equation}
defines a probability measure on $X$ without loss of generality,
since the integral $\int_{X}v(m_{\phi})\omega_{\phi}^{n}$
is independent of the choice of $\phi\in\mathcal{H}^{T}$ \cite{Lah19}.
We denote the ordinary Monge-Amp\`ere measure
$\omega_{\phi}^{n}/\int_{X}\omega_{0}^{n}$
by $\MA(\phi)$.

Following \cite{DJL24}, we introduce generalized notions of
the trace and the Laplacian in the weighted setting.

\begin{definition}(\cite[Appendix A]{DJL24})
For any $T$-invariant $(1,1)$-form $\theta$ with a moment map
$m_{\theta}$, the $v$-weighted trace of $\theta$ is defined by
\begin{equation}
\tr_{v,\phi}(\theta):=\tr_{\phi}(\theta)
+\langle (\log v)^{\prime}(m_{\phi}),m_{\theta}\rangle,
\end{equation}
where $\tr_{\phi}(\theta)$ is the ordinary trace defined by the equation
$\tr_{\phi}(\theta)\omega_{\phi}^{n}=n\theta\wedge\omega_{\phi}^{n-1}$.
Equivalently, in a fixed basis $(\xi_{1},\dots,\xi_{r})$ of $\mathfrak{t}$,
\begin{equation}
\tr_{v,\phi}(\theta)=\tr_{\phi}(\theta)+
\frac{1}{v(m_{\phi})}\sum_{j}v_{, j}(m_{\phi})m^{j}_{\theta},
\end{equation}
where $v_{, j}:=\frac{\partial v}{\partial x_{j}}$ and
$m_{\theta}=(m_{\theta}^{1},\dots,m_{\theta}^{r}).$
\end{definition}

\begin{definition}(\cite[Appendix A]{DJL24})
For any smooth function $f$ on $X$, the $v$-weighted Laplacian
for $\phi\in\mathcal{H}^{T}$ is defined by
\begin{eqnarray}
\Delta_{v,\phi}f&:=&-\frac{1}{v(m_{\phi})}
\overline{\partial}^{*}(v(m_{\phi})\overline{\partial}f)
= \Delta_{\phi} f +\frac{1}{v(m_{\phi})}
\langle\overline{\partial}v(m_{\phi}),\overline{\partial}f\rangle,
\end{eqnarray}
where $\Delta_{\phi}$ is the ordinary (negative) Laplacian for $\phi$.
In particular $\Delta_{v,\phi}f=\tr_{v,\phi}(\deldel f)$.
\end{definition}

It follows from the definition that
the operator $\Delta_{v,\phi}$ is elliptic and 
self-adjoint with respect to
the $v$-twisted Hermitian inner product
$\ip{f_{1}, f_{2}}_{v, \phi}:=\int_{X}f_{1}\overline{f_{2}}\,\MA_{v}(\phi)$
on $C^{\infty}(X;\mathbb{C})$,
and the kernel of $\Delta_{v,\phi}$ consists of the constant functions on $X$.


By considering the weighted Monge-Amp\`ere measure and
the weighted trace, one can define the notions of weighted curvatures.

\begin{definition}
The $v$-weighted Ricci curvature for $\phi\in\mathcal{H}^{T}$ is defined by
the formula
\begin{equation}
\Ric_{v}(\omega_{\phi}):=-\deldel\log\MA_{v}(\phi)
=\Ric(\omega_{\phi})-\deldel\log v(m_{\phi}),
\end{equation}
where $\Ric(\omega_{\phi}):=-\deldel\log\MA(\phi)$
is the ordinary Ricci curvature for $\phi$.
The $v$-weighted scalar curvature for $\phi\in\mathcal{H}^{T}$ is defined by
the formula
\begin{equation}
S_{v}(\phi):=\tr_{v,\phi}\Ric_{v}(\omega_{\phi}).
\end{equation}
\end{definition}

\subsection{Weighted CSCK metrics}
In order to define a weighted CSCK metric,
let $w$ be a smooth function on $P_{X}$ satisfying
$\int_{X}S_{v}(\phi)\MA_{v}(\phi)=\int_{X}w(m_{\phi})\MA_{v}(\phi)$.
Note that the integrals on both sides are independent of the choice
of a metric $\phi\in\mathcal{H}^{T}$ \cite{Lah19}.

\begin{definition}(\cite{Lah19})
A metric $\phi\in\mathcal{H}^{T}$ is called a $(v,w)$ constant scalar curvature
K\"ahler metric (a $(v,w)$-CSCK metric for short) if
\begin{equation}
S_{v}(\phi)=w(m_{\phi}).
\end{equation}
\end{definition}

\begin{remark}\label{w-scalar}
Lahdili \cite[Section 2]{Lah19} originally defined the weighted scalar curvature as
$v(m_{\phi})\tr_{v,\phi}\Ric_{v}(\omega_{\phi})$.
We adopted the definition introduced by Boucksom-Jonsson-Trusiani
\cite[Section 3]{BJT24}.
Namely, a $(v,w)$-CSCK metric as we adopted in this article is equivalent to a $(v,vw)$-CSCK
metric in the sense of Lahdili.
\end{remark}

Typical situations in which a weighted CSCK metric appears are as follows.

\begin{itemize}
\item If $v$ and $w$ are constant function, the weighted CSCK metric
is the ordinary CSCK metric.
\item If $v$ is constant and $w$ is an affine function, the weighted CSCK metric
is the extremal K\"ahler metric introduced by Calabi \cite{Ca82}.
\item If $v$ and $w$ are appropriate polynomials, 
the weighted CSCK metric describes Calabi's extremal metric
on the total space of an holomorphic fibration $Y$ with fiber $X$,
called semisimple principle fibration \cite{ACGT11, AJL23, Jub23}.
\item If $v=e^{l}$ for an affine function and $w=l+a$ for a constant
$a\in\mathbb{R}$ determined by the relation \eqref{w-formula},
the weighted CSCK metric is the $\mu$-CSCK metric,
in the sense of Inoue \cite{Ino22},
extending the notion of K\"ahler-Ricci solitons on Fano manifolds
to general K\"ahler manifolds.
\item If $\Omega=2\pi c_{1}(L)$ for an ample line bundle over $X$, 
$v=l^{-n-1}$ for a positive affine function
and $w=al^{-1}$ for a constant $a\in\mathbb{R}$,
then the weighted CSCK metric describes
a constant scalar curvature Sasaki metric
on the unit circle bundle associated with $(L^{-1})^{\times}$.
In the same setting as above, if one instead takes $w=l^{\prime}l^{-1}$
for an affine function $l^{\prime}$, the weighted CSCK metric
describes a Sasaki extremal metric
on the unit circle bundle associated with $(L^{-1})^{\times}$.
See \cite{AC21, ACL21}.
\item If $v=l^{-2n+1}$ and $w=l^{-2}$ for a positive affine function $l$,
the weighted CSCK metric describes a K\"ahler metric
which is conformal to an Einstein-Maxwell metric \cite{AM19}.
\end{itemize}
%%%%%%%%%%%%%%%%%%%%%%%%%%%%%%%%%%%%
\subsection{Soliton-type metrics}
Let us
assume that $X$ is Fano and $\Omega=2\pi c_{1}(X)$.
For any metric $\phi\in\mathcal{H}^{T}$,
one obtains the Ricci potential function $\rho_{\phi}$ defined by
\begin{equation}
\Ric(\omega_{\phi})-\omega_{\phi}=\deldel\rho_{\phi}
    \quad\text{and}\quad
    \int_{X}(e^{\rho_{\phi}}-1)\MA(\phi)=0.
\end{equation}
In order to introduce soliton-type metrics,
fix a strictly positive and smooth function $g$ on $P_{X}$
such that $\MA_{g}(\phi)$ is a probability measure on $X$
for any $\phi\in\mathcal{H}^{T}$.
One can also define the $g$-weighted Ricci potential function
$\rho_{g,\phi}$ by the formulas
\begin{equation}\label{Ricg}
\Ric_{g}(\omega_{\phi})-\omega_{\phi}=\deldel\rho_{g,\phi}
    \quad\text{and}\quad
    \int_{X}(e^{\rho_{g,\phi}}-1)\MA_{g}(\phi)=0.
\end{equation}
It follows from the definitions that
$\rho_{g,\phi}=\rho_{\phi}-\log g(m_{\phi})$.

\begin{definition}(\cite{BW14, HaLi20})
A metric $\phi\in\mathcal{H}^{T}$ is called a $g$-soliton if
\begin{equation}
\Ric_{g}(\omega_{\phi})=\omega_{\phi},
\quad\text{equivalently,}\quad
\rho_{g,\phi}=0.
\end{equation}
\end{definition}

Let $\mu_{\phi}$ be the canonical measure for $\phi\in\mathcal{H}^{T}$
defined by
\begin{equation}\label{can measure}
\mu_{\phi}=
\frac{e^{-\phi+\rho_{0}}\MA(0)}{\int_{X}e^{-\phi+\rho_{0}}\MA(0)}.
\end{equation}
A $g$-soliton $\phi$ is characterized by
the complex Monge-Amp\`ere equation
$\MA_{g}(\phi)=\mu_{\phi}$.
We will use the following formula frequently later in this paper.
\begin{equation}
\mu_{\phi}=e^{\rho_{\phi}}\MA(\phi)
=e^{\rho_{g,\phi}}\MA_{g}(\phi).
\end{equation}

A $g$-soliton is a weighted CSCK metric in $2\pi c_{1}(X)$
with respect to appropriate weight functions.

\begin{proposition}{\rm (\cite{AJL23})}
A metric $\phi\in\mathcal{H}^{T}$ is a $g$-soliton
if and only if
it is a $(v,w)$-CSCK metric with the weights defined by, for any $x\in P_{X}$,
\begin{equation}\label{w-formula}
v(x)=g(x)\quad\text{and}\quad w(x)=n+\langle(\log g)^{\prime}(x),x\rangle.
\end{equation}
\end{proposition}

\begin{proof}
In general, by taking the $g$-weighted trace of the defining equation
\eqref{Ricg} of the $g$-weighted Ricci potential $\rho_{g,\phi}$,
one has the formula
\begin{equation}\label{LapRicg}
\Delta_{g,\phi}\rho_{g,\phi}
=S_{g}(\phi)-n-\langle(\log g)^{\prime}(m_{\phi}),m_{\phi}\rangle.
\end{equation}
The proposition follows directly from the formula \eqref{LapRicg}.
\end{proof}

Typical situations in which a $g$-soliton appears are as follows.

\begin{itemize}
\item If $g$ is constant, the $g$-soliton is a K\"ahler-Einstein metric.
\item If $g=e^{l}$ for an affine function $l$,
the $g$-soliton is a K\"ahler-Ricci soliton \cite{TZ00}.
\item If $g$ is a positive affine function, the $g$-soliton
is a Mabuchi soliton \cite{Ma01}.
\item If $X$ is toric Fano and $g$ is an appropriate polynomial, the $g$-soliton
describes a K\"ahler-Einstein metric on a total space of fibration $Y$ with fiber $X$,
called KSM manifolds \cite{Nakg19, NN24}.
\item If $g=l^{-n-2}$ for a positive affine function $l$,
then the $g$-soliton describes a Sasaki-Einstein metric on
the unit circle bundle associated with $(K_{X})^{\times}$ \cite{AJL23}.
\end{itemize}
%%%%%%%%%%%%%%%%%%%%%%%%%%%%%%%%%%
\section{$g(v,w)$-solitons}
In this section we explain some background for introducing $g(v,w)$-solitons
from viewpoints of Futaki-type invariants and energy functionals.

\subsection{Futaki-type invariants}
Let $(X, \Omega)$ be a compact K\"ahler manifold.
For any affine function $l$ on $\mathfrak{t}^{*}$, we define
\begin{equation}
M^{\prime}_{v,w}(l)
:=-\int_{X}l(m_0)\Big(S_{v}(0)-w(m_0)\Big)\MA_{v}(0).
\end{equation}
Lahdili \cite{Lah19} showed that
the quantity $M^{\prime}_{v,w}(l)$ is independent of the choice
of a $T$-invariant metric $\omega_{0}\in\Omega$.
Thus if there exists a $T$-invariant $(v,w)$-CSCK metric in $\Omega$
then $M^{\prime}_{v,w}$ must vanish identically.
The invariant $M^{\prime}_{v,w}$ gives an obstruction to the existence
of a $T$-invariant $(v,w)$-CSCK metric.
We call $M^{\prime}_{v,w}$ the $(v,w)$-Futaki invariant.

Let us assume that $X$ is Fano and $\Omega=2\pi c_{1}(X)$.
Let us consider describing the $(v,w)$-Futaki invariant
using the $v$-weighted Ricci potential $\rho_{v}$.
For this we use the following lemma.
\begin{lemma}\label{Deltal}
Fix a holomorphic vector field $\xi\in\mathfrak{t}$
with the potential function $l(m_{\phi})$ for an affine function $l$.
The following equation holds.
\begin{equation}
\Delta_{v,\phi}l(m_{\phi})+l(m_{\phi})+\xi(\rho_{v,\phi})
=\int_{X}l(m_{\phi})e^{\rho_{v,\phi}}\MA_{v}(\phi).
\end{equation}
\end{lemma}
\begin{proof}
Note that $L_{\xi}\omega_{\phi}=\deldel l(m_{\phi})$
by definition of the moment map.
By the Lie derivative of the defining equation
of the weighted Ricci potential \eqref{Ricg},
\begin{equation}
\Delta_{v,\phi}l(m_{\phi})+l(m_{\phi})+\xi(\rho_{v,\phi})=C
\end{equation}
for a constant $C$ which is determined by the normalization 
\begin{equation}
0=\int_{X}L_{\xi}\Big(e^{\rho_{v,\phi}}\MA_{v}(\phi)\Big)
=\int_{X}\Big(\Delta_{v,\phi}l(m_{\phi})+\xi(\rho_{v,\phi})\Big)
e^{\rho_{v,\phi}}\MA_{v}(\phi).
\end{equation}
\end{proof}

Then the weight function
\begin{equation}\label{g(v,w)}
g(v,w)(x):=v(x)\left((1+n+\langle(\log v)^{\prime}(x),x\rangle -w(x)\right)
\end{equation}
naturally appears,
as shown in the following proposition.

\begin{proposition}\label{FutMD}
For any affine function $l$ on $\mathfrak{t}^{*}$, we define
\begin{eqnarray}
D^{\prime}_{g(v,w)}(l)&:=&-\int_{X}l(m_{0})\Big(\frac{g(v,w)(m_{0})}{v(m_{0})}
-e^{\rho_{v,0}}\Big)\MA_{v}(0) \\
&=&-\int_{X}l(m_{0})\Big(g(v,w)(m_{0})-e^{\rho_{0}}\Big)\MA(0).
\end{eqnarray}
Then we have
$M^{\prime}_{v,w}(l)=D^{\prime}_{g(v,w)}(l)$.
\end{proposition}

\begin{proof}
Use the formula \eqref{LapRicg}
and Lemma \ref{Deltal} to show
\begin{eqnarray*}
M^{\prime}_{v,w}(l)&=&\int_{X}
-\langle \dbar l(m_{0}),\dbar\rho_{v,0}\rangle\MA_{v}(0)
+\int_{X}l(m_{0})\Big(\frac{g(v,w)(m_{0})}{v(m_{0})}-1\Big)\MA_{v}(0) \\
&=& \int_{X}-\xi(\rho_{v,0})\MA_{v}(0)
+\int_{X}l(m_{0})\Big(\frac{g(v,w)(m_{0})}{v(m_{0})}-1\Big)\MA_{v}(0) \\
&=&\int_{X}\Big(\Delta_{v,0}l(m_{0})+l(m_{0})
-\int_{X}l(m_{0})e^{\rho_{v,0}}\MA_{v}(0)\Big)\MA_{v}(0) \\
&&+\int_{X}l(m_{0})\Big(\frac{g(v,w)(m_{0})}{v(m_{0})}-1\Big)\MA_{v}(0) \\
&=&D^{\prime}_{g(v,w)}(l).
\end{eqnarray*}
\end{proof}

In other words, the existence of a $(v,w)$-CSCK metric and that of a $g(v,w)$-soliton
are obstructed by the same Futaki-type invariant for a Fano manifold $X$.
Proposition \ref{FutMD} provides evidence for the validity of the statement
that the existence conditions for
a $(v,w)$-CSCK metric and a $g(v,w)$-soliton are equivalent.

\subsection{Energy functionals}
Let $(X,\Omega)$ be a compact K\"ahler manifold.
Analogously to Mabuchi's definition of the Mabuchi functional \cite{Ma86}
by integrating the Futaki invariant \cite{Fu83},
one has a functional on $\mathcal{H}^{T}$
admitting a $(v,w)$-CSCK metric as its critical point
by integrating $M_{v,w}^{\prime}$.
The $(v,w)$-weighted Mabuchi functional
$M_{v,w}:\mathcal{H}^{T}\to\mathbb{R}$ is defined by
\begin{equation}
M_{v,w}(\phi):=-\int_{0}^{1}dt\int_{X}\dot{\phi}_{t}
\Big(S_{v}(\phi_{t})-w(m_{\phi_{t}})\Big)\MA_{v}(\phi_{t}),
\end{equation}
where $\phi_{t}$ is a path in $\mathcal{H}^{T}$ connecting $0$ and $\phi$.
The value $M_{v,w}(\phi)$ is independent of the choice of this path.
Indeed, a so-called Chen-Tian-type formula,
which gives a path-independent expression, is known \cite{BJT24, Lah19}.

Let us assume that $X$ is Fano and $\Omega=2\pi c_{1}(X)$.
Fix a smooth positive function $g$ on $P_{X}$.
Similar to $M_{v,w}$, one has a functional admitting a $g$-soliton
as its critical point.
The $g$-weighted Ding functional $D_{g}:\mathcal{H}^{T}\to\mathbb{R}$
is defined by
\begin{equation}
D_{g}(\phi):=-\int_{0}^{1}dt\int_{X}\dot{\phi}_{t}
\Big(g(m_{\phi_{t}})-e^{\rho_{\phi_{t}}}\Big)\MA(\phi_{t}),
\end{equation}
where $\phi_{t}$ is a path in $\mathcal{H}^{T}$ connecting $0$ and $\phi$.
Put
\begin{equation}\label{E and L}
E_{g}(\phi):=\int_{0}^{1}dt\int_{X}\dot{\phi}_{t}\MA_{g}(\phi_{t})
\quad\text{and}\quad
L(\phi):=-\log\int_{X}e^{\rho_{0}-\phi}\MA(0).
\end{equation}
Then $D_{g}=-E_{g}+L$.
These functionals are defined independently of the choice of a path $\phi_{t}$
\cite{BW14, HaLi20}.

The following proposition,
which plays important roles in our proof of Theorem \ref{main theorem},
may be viewed as an energy functionals version of Proposition \ref{FutMD}.
\begin{proposition}\label{comparison}
For any $\phi\in\mathcal{H}^{T}$, one has
\begin{equation}\label{M=D}
M_{v,w}(\phi)-D_{g(v,w)}(\phi)
=-\int_{X}\rho_{v,\phi}\MA_{v}(\phi)+\int_{X}\rho_{v,0}\MA_{v}(0).
\end{equation}
In particular, 
\begin{equation}\label{M>D}
M_{v,w}(\phi)-D_{g(v,w)}(\phi)\geq \int_{X}\rho_{v,0}\MA_{v}(0).
\end{equation}
\end{proposition}

\begin{proof}
Use the variation formula
$\delta\rho_{v,\phi}=-\Delta_{v,\phi}\delta\phi-\delta\phi+\int_{X}\delta\phi d\mu_{\phi}$,
where $\mu_{\phi}$ is the canonical measure defined as \eqref{can measure},
to show
\begin{eqnarray*}
&&\delta\Big(-\int_{X}\rho_{v,\phi}\MA_{v}(\phi)\Big)\\
&=&\int_{X}\Big(\Delta_{v,\phi}\delta\phi+\delta\phi-\int_{X}\delta\phi d\mu_{\phi}\Big)
\MA_{v}(\phi)
-\int_{X}\rho_{v,\phi}(\Delta_{v,\phi}\delta\phi)\MA_{v}(\phi) \\
&=&\int_{X}\delta\phi\Big[\Big(\frac{g(v,w)(m_{\phi})}{v(m_{\phi})}
+w(m_{\phi})\Big)\MA_{v}(\phi)-d\mu_{\phi}\Big]
-\int_{X}\delta\phi S_{v}(\phi)\MA_{v}(\phi) \\
&=& \delta\left( M_{v,w}(\phi)-D_{g(v,w)}(\phi) \right).
\end{eqnarray*}
By integration, one obtains the equality \eqref{M=D}.
The inequality \eqref{M>D} follows from
the elementary inequality $e^{x}\geq 1+x$ and the normalization of $\rho_{v,\phi}$,
that is,
\begin{equation}
-\int_{X}\rho_{v,\phi}\MA_{v}(\phi)\geq\int_{X}(1-e^{\rho_{v,\phi}})\MA_{v}(\phi)=0.
\end{equation}
\end{proof}
%%%%%%%%%%%%%%%%%%%
\section{Proof of Theorem \ref{main theorem}}
In this section we prove Theorem \ref{main theorem}.
Our proof uses the coercivity theorems which characterizes
the existence of a canonical metric in terms of a behavior of a corresponding
functional on the space of K\"ahler metrics.
Let $I:\mathcal{H}^{T}\to\mathbb{R}$ be the Aubin's $I$-functional defined by
\begin{equation}
I(\phi):=\int_{X}\phi\left(\MA(0)-\MA(\phi)\right),
\end{equation}
which plays a role of a distance function in $\mathcal{H}^{T}$.
In the following theorem, we say that a functional $F:\mathcal{H}^{T}\to\mathbb{R}$ is $T^{\mathbb{C}}$-coercive if
there exists a positive constant $\e>0$ such that the inequality
\begin{equation}
F(\phi)\geq\e\inf_{\sigma\in T^{\mathbb{C}}}I(\sigma[\phi])-\e^{-1}
\end{equation}
holds for any $\phi\in\mathcal{H}^{T}$,
where $\sigma[\phi]\in\mathcal{H}^{T}$ is defined by 
$\sigma^{*}\omega_{\phi}=\omega_{0}+\deldel\sigma[\phi]$.
\begin{theorem}\label{coercivity}
Assume that the weight $v$ is log-concave on $P_{X}$.
\begin{itemize}
\item {\rm (\cite{DJL24, DJL25, HanLiu25})}
There exists a $(v,w)$-CSCK metric in $\mathcal{H}^{T}$ if and only if
the weighted Mabuchi functional $M_{v,w}$ is $T^{\mathbb{C}}$-coercive.
\item {\rm (\cite{HaLi20})} Let $X$ be a Fano manifold and $\Omega=2\pi c_{1}(X)$.
There exists a $v$-soliton in $\mathcal{H}^{T}$ if and only if
the weighted Ding functional $D_{v}$ is $T^{\mathbb{C}}$-coercive.
\end{itemize}
\end{theorem}

Let us focus on the case where $X$ is Fano and $\Omega=2\pi c_{1}(X)$
to prove Theorem \ref{main theorem}.
One direction is easier to follow.
In fact it follows from Theorem \ref{coercivity} and
the inequality \eqref{M>D} in Proposition \ref{comparison}
that the existence of a $g(v,w)$-soliton implies that of a $(v,w)$-CSCK metric.
As for the converse direction, there is currently no known method 
to directly deduce the coercivity of $D_{g(v,w)}$ from that of $M_{v,w}$.

Our approach is to construct a $g(v,w)$-soliton 
by a continuity method under the assumption on
the $T^{\mathbb{C}}$-coercivity of $M_{v,w}$.
Let us consider the following PDE for $\phi\in\mathcal{H}^{T}$:
\begin{equation}\label{conti-eq}
\MA_{g(v,w)}(\phi)=e^{-t\phi+\rho_{0}}\MA(0),
\end{equation}
where $t\in [0,1]$.
A solution of the equation \eqref{conti-eq} at $t=1$ is nothing but a $g(v,w)$-soliton.
Let $\mathcal{T}$ be the set of $t\in[0,1]$ such that the equation \eqref{conti-eq}
has a solution in $\mathcal{H}^T$.
According to \cite[Section 3]{HN24},
the set $\mathcal{T}$ is non-empty and open in general.
Moreover, according to \cite[Section 5]{HN24},
if for any $\e\in(0,1)$ there exists a constant $C>0$ independent of $t$ such that
\begin{equation}
\inf_{\sigma\in T^{\mathbb{C}}}I(\sigma[\phi_{t}])\leq C
\end{equation}
holds for any solution $\phi_{t}\in\mathcal{H}^{T}$ of \eqref{conti-eq}
at any $t\in\mathcal{T}\cap(\e,1)$,
then the set $\mathcal{T}$ is closed, which implies $1\in\mathcal{T}$.
In view of the assumption on the $T^{\mathbb{C}}$-coercivity of $M_{v,w}$,
the rest of our proof of Theorem \ref{main theorem} is deduced to
that of the following.

\begin{theorem}\label{bound of M}
Fix any $\e\in(0,1)$.
Let $\phi_t$ be a $T$-invariant solution of \eqref{conti-eq} at $t\in(\e,1)$.
There exists a constant $C>0$ independent of $\e$ and $t$ satisfying
\begin{equation}\label{estimate of M}
M_{v,w}(\phi_t)\leq C\e^{-1}.
\end{equation}
\end{theorem}

In the rest of this section we prove Theorem \ref{bound of M}.
Let $\phi_{t}\in\mathcal{H}^T$ be a solution of the equation \eqref{conti-eq} at $t$.
In order to obtain the upper bound of $M_{v,w}(\phi_{t})$, 
it suffices to control $D_{g(v,w)}(\phi_{t}) -\int_{X}\rho_{v,\phi_{t}}\MA_{v}(\phi_{t})$
from above by the formula \eqref{M=D} in Proposition \ref{comparison}.
Put $\omega_{\phi_{t}}:=\omega_{t}$,
$\rho_{t}:=\rho_{\phi_{t}}$, $\rho_{v,t}:=\rho_{v,\phi_{t}}$,
$v_{t}:=v(m_{\phi_{t}})$ and $g_{t}:=g(v,w)(m_{\phi_{t}})$ for simplicity.

By definition of the weighted Ding functional $D_{g(v,w)}$ and \cite[Proposition 4.4]{HN24},
\begin{equation}\label{Dg-estimate}
D_{g(v,w)}(\phi_{t})=-E_{g(v,w)}(\phi_{t})+L(\phi_{t})\leq L(\phi_{t}).
\end{equation}
On the other hand, we observe the following to control 
$-\int_{X}\rho_{v,\phi_{t}}\MA_{v}(\phi_{t})$.

\begin{lemma}\label{rho-eq}
Let $\phi_{t}\in \mathcal{H}^T$ be a solution of \eqref{conti-eq} at $t$. 
One has the formula
\begin{equation}\label{rho-formula}
\rho_{v,t}=-(1-t)\phi_{t}+\log\frac{g_{t}}{v_{t}}+L(\phi_{t}),
\end{equation}
where $L$ is the functional defined in \eqref{E and L}.
\end{lemma}
Note that by the compactness of $X$,
there exists a uniform constant $C>0$ independent of
the choice of $\phi\in\mathcal{H}^{T}$ such that
\begin{equation}
C^{-1}\leq \frac{g(v,w)(m_{\phi})}{v(m_{\phi})} \leq C.
\end{equation}

\begin{proof}
By definition of the weighted Ricci curvature,
\begin{equation}\label{wRic-formula}
\Ric_{v}\omega_{t}=\deldel\log\frac{g_{t}}{v_{t}}+\Ric_{g(v,w)}(\omega_{t}).
\end{equation}
By the formula \eqref{wRic-formula}, definition of the $v$-weighted Ricci potential 
and the equation \eqref{conti-eq},
\begin{equation}
\deldel\rho_{v,t}=\deldel\left(\log\frac{g_{t}}{v_{t}}-(1-t)\phi_{t}\right).
\end{equation}
Thus $\rho_{v,t}=-(1-t)\phi_{t}+\log g_{t}/v_{t}+c_{t}$
for some constant $c_{t}$ depending on $t$.
The constant $c_{t}$ is determined by the normalization of $\rho_{v,t}$
and the equation \eqref{conti-eq}.
In fact $c_{t}=L(\phi_{t})$.
This completes the proof.
\end{proof}

Integrating the both side of the relation in Lemma \ref{rho-eq},
one thus obtains
\begin{equation}
-\int_{X}\rho_{v,\phi_{t}}\MA_{v}(\phi_{t})
\leq
(1-t)\int_{X}\phi_{t}\MA_{v}(\phi_{t}) -\log\left(\inf_{P_{X}} \frac{g(v,w)}{v}\right)-L(\phi_t).
\end{equation}

\begin{lemma}\label{int-phi}
For a fixed $\e\in (0,1)$,
a solution $\phi_{t}$ at $t>\e$ satisfies
\begin{equation}
\int_{X}\phi_{t}\MA_{v}(\phi_{t})
\leq e^{\sup\rho_{0}-1}\left(\inf_{P_{X}} \frac{g(v,w)}{v}\right)^{-1}\e^{-1}.
\end{equation}
\end{lemma}

\begin{proof}
Using the equation \eqref{conti-eq} of the continuity method,
one obtains
\begin{eqnarray*}
\int_{X}\phi_{t}\,\MA_{v}(\phi_{t})
&=&
\int_{X}\phi_{t}\frac{v_{t}}{g_{t}}\,\MA_{g}(\phi_{t})\\
&\leq&
\int_{\{\phi_{t}>0\}}\phi_{t}\frac{v_{t}}{g_{t}}e^{-t\phi_{t}+\rho_{0}}
\,\MA(0)\\
&\leq&
e^{\sup\rho_{0}}\left(\inf_{P_{X}} \frac{g(v,w)}{v}\right)^{-1}
\int_{\{\phi_{t}>0\}}\phi_{t}e^{-t\phi_{t}}\MA(0).
\end{eqnarray*}
Note that for any $t>\e$, the function $\mathbb{R}_{>0}\ni x\mapsto xe^{-tx}$
is bounded from above by the constant $(e\e)^{-1}$.
This completes the proof.
\end{proof}

Putting together the above arguments,
Theorem \ref{bound of M} follows
from the relation \eqref{M=D}, the estimate \eqref{Dg-estimate} and Lemma \ref{int-phi}.
This completes the proof of Theorem \ref{main theorem}.
%%%%%%%%%%%%%%%%%%%%%%%%%%%%%%%%%%%%%
\section{$(v,w)$-CSCK metrics and $g(v,w)$-solitons in Sasaki geometry}
\label{Sasaki section}
The purpose of this section is to prove Theorem \ref{Sasaki theorem}
(and Theorem \ref{Sasaki g(v,w)} for more general viewpoint)
explaining how the weight function $g(v,w)$
arises naturally from Sasaki geometry.
For this, we discuss Sasaki-extremal metrics and Sasaki-Mabuchi solitons
(for a possibly irregular Reeb field)
on a Sasaki manifold admitting a regular Reeb field.
We shall see
that a Sasaki-extremal metric and a Sasaki-Mabuchi soliton give rise,
via the regular quotient, to a $(v,w)$-CSCK metric and a $g(v,w)$-soliton
respectively on the
underlying Fano manifold with the same weights $(v,w)$.
More generally our argument also applies to a transverse $g$-extremal metric
and a transverse $g$-soliton. 

We refer to \cite{BGBook, MSY08} for basics of Sasaki manifolds.
We also refer to \cite{FOW09, CL21} for differential geometry in Sasaki manifolds
including transverse K\"ahler geometry.

\subsection{Basic setup}
Let $(N,g)$ be a compact Sasaki manifold of real dimension $2n+1$ for $n\geq 1$ i.e.
the cone $(C(N), \bar{g}):=(N\times\mathbb{R}_{+}, r^{2}g+dr^{2})$ is a K\"ahler manifold
of complex dimension $n+1$.
The manifold $N$ is often identified with the submanifold $\{r=1\}\subset C(N)$.
A Sasaki manifold $(N,g)$ admits a contact structure $(\hat{\xi}, \eta,\Phi)$
and also admits a one dimensional foliation $\mathcal{F}_{\hat{\xi}}$.
Here the killing vector field $\hat{\xi}$ on $(N,g)$ is called the Reeb vector field,
$\eta$ is called the contact $1$-form $N$,
$\Phi$ is a $(1,1)$ tensor on $N$ which defines a complex structure on the contact sub-bundle
$D:=\Ker\eta$.
A Sasaki manifold is also denoted by $(N,\hat{\xi},\eta,D,\Phi)$.
The data $(\hat{\xi},\eta,D,\Phi)$ is called a Sasaki structure on $N$.
A Sasaki manifold $(N,\hat{\xi},\eta,D,\Phi)$ is said to be
(i) regular if $\mathcal{F}_{\hat{\xi}}$ is obtained by a free $S^{1}$-action,
(ii) quasi-regular if all the leaves of
$\mathcal{F}_{\hat{\xi}}$ is compact
and (iii) irregular if it is not quasi-regular.

The Sasaki structure $(\hat{\xi},\eta,D,\Phi)$ on $N$
yields a transverse holomorphic structure
and a transverse K\"ahler structure on $\mathcal{F}_{\hat{\xi}}$.
One has the splitting of
complex valued basic $p$-forms and the derivative:
\begin{equation}
\Lambda_{B}(N)^{p}=\oplus_{i+j=p}\Lambda_{B}^{i,j}(N),
\quad d_{B}:=d|_{\Lambda_{B}^{p}(N)}=\partial_{B}+\dbar_{B}.
\end{equation}
Put $d_{B}^{c}=\frac{\sqrt{-1}}{2}(\dbar_{B}-\partial_{B})$ so that
$d_{B}d_{B}^{c}=\deldelb$.
If we emphasize the Reeb field $\hat{\xi}$, we write
$d_{\hat{\xi}}, \partial_{\hat{\xi}}, \dbar_{\hat{\xi}}$ and $d^{c}_{\hat{\xi}}$.
Let $\omega$ 
be a basic K\"ahler form on $N$.
The transverse Ricci form $\Ric^{T}$ for $\omega$ is defined as
$\Ric^{T}=-\deldelb\log\det\omega^{n}$.
The basic cohomology class represented by $\Ric^{T}/2\pi$ is independent of the choice
of $\omega$.
This basic cohomology class $[\Ric^{T}/2\pi]_{B,\hat{\xi}}$ is called the basic first Chern class
and is denoted by $c^{B,\hat{\xi}}_{1}(N)$.
%%%%%%%%%%%%%%%%%%%%
\subsection{Transverse $g$-solitons and $g$-extremal metrics}
Let $\Ric^{T}$ is the transverse Ricci form for
$\omega^{T}:=d\eta$.
The transverse scalar curvature $S^{T}$ for $\omega^{T}$ is defined by
$S^{T}(\omega^{T})^{n}\wedge\eta
=n\Ric^{T}\wedge(\omega^{T})^{n-1}\wedge\eta$.
Let $V:=\int_{N}(\omega^{T})^{n}\wedge\eta$ be the volume of $N$
which is independent of the choice of the transverse K\"ahler metric in
$[d\eta]_{B,\hat{\xi}}$.
The average
$\underline{S}^{T}=\int_{N}S^{T}(\omega^{T})^{n}\wedge\eta/V$
of $S^{T}$ is also independent of the choice of $\eta$.

Let $\hat{T}$ be a real maximal compact torus acting effectively as CR-automorphisms of
$(D=\ker\eta, \Phi|_{D})$ such that
$\hat{\xi}\in\hat{\frak{t}}:=\mathrm{Lie}(\hat{T})$.
One has the contact moment map
$\mu^{\hat{\xi}}:C(N)\to\hat{\mathfrak{t}}^{*}$ defined by
\begin{equation}
\langle\mu^{\hat{\xi}},u\rangle=r^{2}\eta(u).
\end{equation}
Let $\mathcal{C}^{*}\subset\hat{\frak{t}^{*}}$ be
the image of $\mu^{\hat{\xi}}$
which is a convex rational polyhedral cone.
The image of the restriction $\mu^{\hat{\xi}}|_{N}$ is given by
\begin{equation}
P_{\hat{\xi}}:=\Set{x\in\mathcal{C}^{*} | \langle x,\hat{\xi}\rangle=1}
\end{equation}
and is called the characteristic hyperplane.
Let $\mathcal{C}:=\Set{y\in\hat{\frak{t}} | \langle x,y\rangle >0 \text{ for any } x\in\mathcal{C}^{*}}$
be the dual cone of $\mathcal{C}^{*}$.
An element $\hat{\chi}\in\hat{\frak{t}}$ lies in $\mathcal{C}$
if and only if $\eta(\hat{\chi})>0$.
Thus $\hat{\chi}\in\mathcal{C}$ defines a Reeb filed on $N$  \cite{HS16}.
The cone $\mathcal{C}$ is called the Reeb cone or the Sasaki cone.
For any Reeb field $\hat{\chi}\in\mathcal{C}$,
there exists a natural projection
\begin{equation}\label{projection}
P_{\hat{\xi}}\ni x\mapsto \frac{x}{\langle x,\hat{\chi}\rangle} \in P_{\hat{\chi}}.
\end{equation}

In order to define the notions of a transverse $g$-extremal metric
and a transverse $g$-soliton,
let $g$ be a smooth positive function on $P_{\hat{\xi}}$
satisfying $\int_{N}g(\mu^{\hat{\xi}})(\omega^{T})^{n}\wedge\eta=V$.
The positivity of the weight function 
$g$ is not necessarily required to define transverse
$g$-extremal metrics, whereas it is essential for defining transverse 
$g$-solitons.
\begin{definition}
A transverse K\"ahler metric $\omega^{T}$ is called a transverse
$g$-extremal metric if
\begin{equation}
S^{T}-\underline{S}^{T}=1-g(\mu^{\hat{\xi}}|_{N}).
\end{equation}
\end{definition}

If $g$ is an affine function, the transverse $g$-extremal metric
is an extremal Sasaki metric in the sense of
Boyer-Galicki-Simanca \cite{BGS08}.

Let us now assume transversely Fano condition
$2\pi c^{B,\hat{\xi}}_{1}(N)=[d\eta]_{B,\hat{\xi}}$.
In this case $\underline{S}^{T}=n$.
The condition $2\pi c^{B,\hat{\xi}}_{1}(N)=[d\eta]_{B,\hat{\xi}}$ follows from
the topological assumptions
$c^{B,\hat{\xi}}_{1}(N)>0$ and $c_{1}(D)=0$
and so-called the $D$-homothetic transformation for the radial function on
the Riemannian cone $C(N)$.
In fact $c_{1}(D)=0$ if and only if $2\pi c^{B,\hat{\xi}}_{1}(N)$ is represented by
$\tau d\eta$ for some constant $\tau$ \cite{FOW09}.

\begin{definition}
A transverse K\"ahler metric $\omega^{T}$ is called a transverse
$g$-soliton if
\begin{equation}
\Ric^{T}-\omega^{T}=\deldelb \log g(\mu^{\hat{\xi}}|_{N})
\end{equation}
\end{definition}

If $g=e^{l}$ for an affine function $l$, the transverse $g$-soliton
is a Sasaki Ricci soliton in the sense of Futaki-Ono-Wang \cite{FOW09}.
If $g$ is an affine function, the transverse $g$-soliton
should be called a Sasaki Mabuchi soliton.
%%%%%%%%%%%%%%%%%%%%
\subsection{Regular quotient}
Let $X$ be a Fano manifold, $T$ a maximal compact torus
in the automorphism group $\mathrm{Aut}(X)$
and $\omega_{0}\in 2\pi c_{1}(X)$ a $T$-invariant K\"ahler metric.
One has the associated moment map $m_{0}:X\to\frak{t}^{*}$ with
the moment polytope $P_{X}$.
Let us consider the unitary $S^{1}$-bundle 
$\pi : N_{\omega_{0}}\to X$ contained in the canonical bundle $K_{X}$
with respect to the Hermitian metric on $K_{X}$
whose curvature is $-\omega_{0}$.
The new torus $\hat{T}:=T\times S^{1}$ acts on $N_{\omega_{0}}$,
where $S^{1}$ acts on the fibers,
preserving the induced CR structure $(D_{0}, J_{0})$.
Here $D_{0}$ is the horizontal distribution
and $J_{0}$ is induced complex structure on $D_{0}$.
In fact $\hat{T}$ is a maximal compact torus acting effectively
as CR automorphisms.
Let $\hat{\chi}\in\hat{\frak{t}}:=\mathrm{Lie}(\hat{T})$
be the generator of the $S^{1}$-action.
The corresponding $1$-form $\eta_{0}^{\hat{\chi}}$ is uniquely determined
as the connection $1$-form on $N_{\omega_{0}}$ with the curvature form
$d\eta_{0}^{\hat{\chi}}=\pi^{*}\omega_{0}$.
The kernel of $\eta^{\hat{\chi}}$ is nothing but the horizontal distribution $D_{0}$.
Thus the transverse K\"ahler structure for the Sasaki structure
$(\hat{\chi}, \eta_{0}^{\hat{\chi}}, D_{0},J_{0})$ on $N_{\omega_{0}}$
is the pull back of $\omega_{0}$ to $(D_{0},J_{0})$.
In particular the transverse Fano condition
$2\pi c^{B,\hat{\chi}}_{1}(N_{\omega_{0}})
=[d\eta_{0}^{\hat{\chi}}]_{B,\hat{\chi}}$ holds.

Let $\xi\in\frak{t}$ to define the affine linear function
$l_{\xi}(x)=\langle x,\xi \rangle +1$ on $P_{X}$.
Let us consider
\begin{equation}
\hat{\xi}:=\xi+\hat{\chi}\in\hat{\frak{t}}.
\end{equation}
An element $\xi\in\frak{t}\subset\hat{\frak{t}}$ gives rise to a vector field on $N_{\omega_{0}}$
given by the horizontal lift of $\xi$ to $D_{0}$ plus
$\langle \pi^{*}m_{0},\xi\rangle\hat{\chi}$.
It follows that $\eta_{0}^{\hat{\chi}}(\hat{\xi})=l_{\xi}(\pi^{*}m_{0})$.
Thus $\hat{\xi}$ defines a Reeb field of $(N_{\omega_{0}},D_{0},J_{0})$
if and only if $l_{\xi}$ is a positive function on $P_{X}$.
Moreover, if $\hat{\xi}$ defines a Reeb field
then the corresponding contact form $\eta_{0}^{\hat{\xi}}$
satisfies the transverse Fano condition
$2\pi c^{B,\hat{\xi}}_{1}(N)=[d\eta_{0}^{\hat{\xi}}]_{B,\hat{\xi}}$
as explained in \cite[Section 4.6]{ALL25} for example.

Assume that $\hat{\xi}=\xi+\hat{\chi}$ defines a Reeb field
of $(N_{\omega_{0}},D_{0},J_{0})$.
Following \cite[Section 2.3]{ACL21}
(building on an idea in \cite[Proof of Lemma 2.2]{HS16}),
we introduce the variation space of Sasaki structures as follows.
\begin{equation}
\Xi_{\hat{\xi},\eta^{\hat{\xi}}_{0},J^{\hat{\xi}}}^{\hat{T}}(N_{\omega_{0}})
:=
\left\{
\varphi\in C^{\infty}(N_{\omega_{0}};\mathbb{R})^{\hat{T}}
\ \middle|\
\eta^{\hat{\xi}}_{\varphi}
:=
\eta^{\hat{\xi}}_{0}+d^{c}_{\hat{\xi}}\varphi
\ \text{satisfies } 
d\eta^{\hat{\xi}}_{\varphi}>0
\ \text{on }D_{\varphi}:=\ker(\eta^{\hat{\xi}}_{\varphi})
\right\},
\tag{KRS11}
\end{equation}
where $J^{\hat{\xi}}\in\mathrm{End}(TN_{\omega_{0}})$ extends
$J_{0}\in\mathrm{End}(D_{0})$ by assuming the conditions
$J^{\hat{\xi}}(\hat{\xi})=0$,
$
d^{c}_{\hat{\xi}}\varphi := -\,d\varphi\circ J^{\hat{\xi}}
$
and that the positivity of $d\eta^{\hat{\xi}}_{\varphi}$ on $D_{\varphi}$
is defined with respect to $J^{\hat{\xi}}$.
We call an element in
$\Xi_{\hat{\xi},\eta^{\hat{\xi}}_{0},J^{\hat{\xi}}}^{\hat{T}}(N_{\omega_{0}})$
a $\hat{\xi}$-transverse K\"ahler potential.
As explained in \cite[Section 2.3]{ACL21},
there is a bijection $\Theta_{\hat{\chi},\hat{\xi}}:
\Xi_{\hat{\xi},\eta^{\hat{\xi}}_{0},J^{\hat{\xi}}}^{\hat{T}}(N_{\omega_{0}})
\to
\Xi_{\hat{\chi},\eta^{\hat{\chi}}_{0},J^{\hat{\chi}}}^{\hat{T}}(N_{\omega_{0}})$,
and $\phi$ and $\Theta_{\hat{\chi},\hat{\xi}}(\phi)$ induce Sasaki structures
with the same underlying CR structure.
Note that 
$\Xi_{\hat{\chi},\eta^{\hat{\chi}}_{0},J^{\hat{\chi}}}^{\hat{T}}(N_{\omega_{0}})$
is identified with $\mathcal{H}^{T}$ via the pullback by $\pi:N_{\omega_{0}}\to X$. 

Now let's see
how a $\hat{\xi}$-transverse $g$-soliton and a $\hat{\xi}$-transverse
$g$-extremal metric are translated into
K\"ahler structures associated with $\hat{\chi}$, namely, 
into the K\"ahler structure of $X$.
For a $\hat{\xi}$-transverse K\"ahler potential $\hat{\phi}$,
one obtains the unique K\"ahler potential $\phi\in\mathcal{H}^{T}$ on $X$
satisfying $\Theta_{\hat{\chi},\hat{\xi}}(\hat{\phi})=\pi^{*}\phi$.
It follows from the formula 
\begin{equation}\label{eta relation}
\eta^{\hat{\xi}}_{\hat{\phi}}
=(\eta^{\hat{\chi}}_{\pi^{*}\phi}(\hat{\xi}))^{-1}\eta^{\hat{\chi}}_{\pi^{*}\phi}
\end{equation}
that the contact moment map $\mu^{\hat{\xi}}_{\hat{\phi}}$
defined by $\eta^{\hat{\xi}}_{\hat{\phi}}$,
the contact moment map $\mu^{\hat{\chi}}_{\pi^{*\phi}}$
defined by $\eta^{\hat{\chi}}_{\pi^{*}\phi}$
and the moment map $m_{\phi}$ defined by $\phi\in\mathcal{H}^{T}$
satisfy the relation
\begin{equation}
\mu^{\hat{\xi}}_{\hat{\phi}}
=\frac{\mu^{\hat{\chi}}_{\pi^{*\phi}}}{\eta^{\hat{\chi}}_{\pi^{*}\phi}(\hat{\xi})}
=\frac{\pi^{*}m_{\phi}+\hat{\chi}^{*}}{l_{\xi}(\pi^{*}m_{\phi})},
\end{equation}
where $\hat{\chi}^{*}\in\hat{\frak{t}}^{*}$ stands as the dual element of
$\hat{\chi}\in\hat{\frak{t}}$.
See \cite[Section 2.4]{CL21} for example.
Thus one obtains a smooth positive function $g_{0}$ on $P_{X}$ satisfying
\begin{equation}
g_{0}(\pi^{*}m_{\phi})=g(\mu^{\hat{\xi}}_{\hat{\phi}})
\end{equation}
on $N_{\omega_{0}}$,
where $g$ is the weight function on $P_{\hat{\xi}}$
in the defining equation of a transverse $g$-extremal metric
and a transverse $g$-soliton. 
This terminology is used in the following propositions,
which show how a $\hat{\xi}$-transverse $g$-soliton and a $\hat{\xi}$-
transverse $g$-extremal metric on $N_{\omega_{0}}$
are translated into a weighted CSCK metric and a soliton-type metric
on $X$ respectively.

\begin{proposition}\label{reg-sigma}
A $\hat{\xi}$-transverse K\"ahler potential
$\hat{\phi}\in
\Xi_{\hat{\xi},\eta^{\hat{\xi}}_{0},J^{\hat{\xi}}}^{\hat{T}}(N_{\omega_{0}})$
defines a $\hat{\xi}$-transverse $g$-soliton
$\omega_{\hat{\xi}}^{T}:=d\eta^{\hat{\xi}}_{\hat{\phi}}$
in $2\pi c^{B,\hat{\xi}}_{1}(N_{\omega_{0}})$
if and only if
the K\"ahler potential $\phi\in\mathcal{H}^{T}$
given by $\Theta_{\hat{\chi},\hat{\xi}}(\hat{\phi})=\pi^{*}\phi$
defines
a $g_{\xi}$-soliton $\omega_{\phi}$ with the weight
\begin{equation}
g_{\xi}=l_{\xi}^{-n-2}g_{0}.
\end{equation}
\end{proposition}

\begin{proof}
We follow an argument in \cite{ALL25} (see also \cite{AJL23, CL21}).
It follows from \cite[Lemma 4.3]{ALL25} and the formula
\eqref{eta relation} that
$\omega_{\hat{\xi}}^{T}$ is a $\hat{\xi}$-transverse $g$-soliton
if and only if the $\hat{\chi}$-transverse K\"ahler metric
$\pi^{*}\omega_{\phi}=d\eta_{\pi^{*}\phi}^{\hat{\chi}}$ satisfies
\begin{equation}
\Ric^{T}(\pi^{*}\omega_{\phi})-\pi^{*}\omega_{\phi}=
dd^{c}_{\hat{\chi}}g_{0}(\pi^{*}m_{\phi})
-(n+2)dd^{c}_{\hat{\chi}}\log l_{\xi}(\pi^{*}m_{\phi}).
\end{equation}
Thus the K\"ahler metric $\omega_{\phi}$ is a $g_{\xi}$-soliton on $X$.
\end{proof}

\begin{proposition}\label{reg-ext}
A $\hat{\xi}$-transverse K\"ahler potential
$\hat{\phi}\in
\Xi_{\hat{\xi},\eta^{\hat{\xi}}_{0},J^{\hat{\xi}}}^{\hat{T}}(N_{\omega_{0}})$
defines a $\hat{\xi}$-transverse $g$-extremal metric
$\omega_{\hat{\xi}}^{T}:=d\eta^{\hat{\xi}}_{\hat{\phi}}$
in $2\pi c^{B,\hat{\xi}}_{1}(N_{\omega_{0}})$
if and only if
the K\"ahler potential $\phi\in\mathcal{H}^{T}$
given by $\Theta_{\hat{\chi},\hat{\xi}}(\hat{\phi})=\pi^{*}\phi$
defines
a $(v_{\xi},w_{\xi})$-CSCK metric with the weights
\begin{equation}
v_{\xi}=l_{\xi}^{-n-1}
\quad\text{and}\quad
w_{\xi}=\left(n+1-g_{0}\right)l_{\xi}^{-1}.
\end{equation}
\end{proposition}

\begin{proof}
We apply the result of \cite[Lemma 3]{AC21} proving that
\begin{equation}
S^{T}(\hat{\phi})\hat{\xi}=
l_{\xi}(\pi^{*}m_{\phi})^{2}(\pi^{*}S_{v_{\xi}}(\phi))\hat{\chi},
\end{equation}
where $S^{T}(\hat{\phi})$ is the transverse scalar curvature of
$\omega^{T}_{\hat{\xi}}$
and $S_{v_{\xi}}$ is the $v_{\xi}$-weighted scalar curvature of
$\omega_{\phi}$.
Recall Remark \ref{w-scalar} for the notation.
It follows from the formula \eqref{eta relation}
that $S^{T}(\hat{\phi})=l_{\xi}(\pi^{*}m_{\phi})\pi^{*}S_{v_{\xi}}(\phi)$.
Thus the defining equation of a $\hat{\xi}$-transverse $g$-extremal metric
is equivalent to
\begin{equation}
\pi^{*}S_{v_{\xi}}(\phi)=w_{\xi}(\pi^{*}m_{\phi}),
\end{equation}
which shows that $\omega_{\phi}$ is a $(v_{\xi},w_{\xi})$-CSCK metric
in $2\pi c_{1}(X)$.
\end{proof}

Summarizing the above arguments,
we have the following theorem, showing that
a $g(v,w)$-soliton appears naturally in Sasaki geometry of $N_{\omega_{0}}$.

\begin{theorem}\label{Sasaki g(v,w)}
The weight $g(v_{\xi},w_{\xi})$ associated with the weight
$(v_{\xi},w_{\xi})$ in Proposition \ref{reg-ext} equals
the wight $g_{\xi}$ in Proposition
\ref{reg-sigma},
that is,
for any $x\in P_{X}$,
\begin{equation}
g(v_{\xi},w_{\xi})(x)=g_{\xi}(x).
\end{equation}
\end{theorem}
\begin{proof}
By definition of $g(v,w)$-weight \eqref{g(v,w)}, one obtains
\begin{eqnarray*}
g(v_{\xi},w_{\xi})(x)
&=&
l_{\xi}^{-n-1}\left( 1+n -(n+1)\langle (\log l_{\xi})^{\prime}(x),x\rangle
-\left(n+1-g_{0}\right)l_{\xi}^{-1}\right) \\
&=&
l_{\xi}^{-n-1}\left( 1+n -(n+1)\frac{l_{\xi}(x)-1}{l_{\xi}}
-\left(n+1-g_{0}\right)l_{\xi}^{-1}\right) \\
&=&g_{\xi}(x).
\end{eqnarray*}
\end{proof}

Theorem \ref{Sasaki theorem} stated in the introduction corresponds to Theorem
\ref{Sasaki g(v,w)}, where the weight function $g$ defining a $\hat{\xi}$-transverse
$g$-extremal metric and a $\hat{\xi}$-transverse $g$-soliton
is assumed to be affine.
%%%%%%%%%%%%%%%%%%%%%%%%%%%%%%%%%%%%%%
\begin{thebibliography}{99}
\bibitem{AC21}
\textsc{V.~Apostolov and D.~M.~J.~Calderbank},
The CR geometry of weighted extremal K\"ahler and Sasaki metrics,
Math. Ann., 379(3-4):1047--1088, 2021.

\bibitem{ACL21}
\textsc{V.~Apostolov, D.~M.~J.~Calderbank and E.~Legendre},
Weighted K-stability of polarized varieties and extremality of Sasaki manifolds,
Adv. Math. 391 (2021), Paper No. 107969, 63 pp.

\bibitem{ACGT11}
\textsc{V.~Apostolov, D.~M.~J.~Calderbank, P.~Gauduchon, and C.~T\o nnesen-Friedman},
Extremal K\"ahler metrics on projective bundles over a curve,
Adv. Math., 227:2385--2424, 2011.

\bibitem{AJL23}
\textsc{V.~Apostolov, S.~Jubert and A.~Lahdili},
Weighted K-stability and coercivity with applications to
extremal K\"ahler and Sasaki metrics,
Geom. Topol. 27 (2023), 3229--3302.

\bibitem{ALL25}
\textsc{V.~Apostolov, A.~Lahdili and E.~Legendre},
From K\"ahler Ricci solitons to Calabi-Yau K\"ahler cones,
Real and complex geometry--in honour of Paul Gauduchon,
Springer, Cham, [2025], 1--40.

\bibitem{ALN25}
\textsc{V.~Apostolov, A.~Lahdili and Y.~Nitta},
Mabuchi K\"ahler solitons versus extremal K\"ahler metrics and beyond,
Bull. Lond. Math. Soc. 57 (2025), no. 3, 692--710.

\bibitem{ALN25b}
\textsc{V.~Apostolov, A.~Lahdili and Y.~Nitta},
Private communication.

\bibitem{AM19}
\textsc{V.~Apostolov and G.~Maschler},
Conformally K\"ahler, Einstein-Maxwell geometry,
J. Eur. Math. Soc. (JEMS), 21 (2019), no. 5, 1319--1360.

\bibitem{BW14}
\textsc{R. Berman, D. Witt Nystr\"{o}m}, 
Complex optimal transport and the pluripotential theory of
K\"{a}hler-Ricci solitons,
arXiv:1401.8264.

\bibitem{BJ25}
\textsc{S.~Boucksom and M.~Jonsson},
On the Yau-Tian-Donaldson conjecture for weighted cscK metrics,
arXiv:2509.15016.

\bibitem{BJT24}
\textsc{S.~Boucksom, M.~Jonsson and A.~Trusiani},
Weighted extremal K\"ahler metrics on resolution of singularities,
Comm. Pure Appl. Math.. https://doi.org/10.1002/cpa.70026

\bibitem{BGBook}
\textsc{C.~P.~Boyer and K.~Galicki}, 
Sasaki Geometry,
Oxford Mathematical Monographs, 2008.

\bibitem{BGS08}
\textsc{C.~P.~Boyer, K.~Galicki and S.~R.~Simanca},
Canonical Sasakian metrics,
Comm. Math. Phys. 279 (2008), no. 3, 705--733.

\bibitem{Ca82}
\textsc{E.~Calabi},
Extremal K\"ahler metrics, Seminar on Differential Geometry
(1982), Vol.102, Princeton University Press.

\bibitem{DJL24}
\textsc{E.~Di~Nezza, S.~Jubert and A.~Lahdili},
Weighted CSCK metric (I): A priori estimates,
J. Funct. Anal. 289 (2025), no. 11, Paper No. 111148, 37 pp.

\bibitem{DJL25}
\textsc{E.~Di~Nezza, S.~Jubert and A.~Lahdili},
Weighted cscK metric (II): the continuity method,
arXiv:2503.22183

\bibitem{Fu83}
\textsc{A.~Futaki},
An obstruction to the existence of Einstein K\"ahler metrics.
Invent. Math. 73 (1983), no. 3, 437--443.

\bibitem{FM95}
\textsc{A. Futaki and T. Mabuchi}, 
Bilinear forms and extremal K\"{a}hler vector fields associated with
K\"{a}hler classes,
Math. Ann. 301(1995), 199--210.

\bibitem{FOW09}
\textsc{A.~Futaki, H.~Ono and G.~Wang},
Transverse K\"ahler geometry of Sasaki manifolds
and toric Sasaki-Einstein manifolds,
J. Differential Geom. 83, (2009) 585--635.

\bibitem{HaLi20}
\textsc{J.~Han and C.~Li},
On the Yau-Tian-Donaldson conjecture for generalized K\"{a}hler-Ricci
soliton equations,
Comm. Pure Appl. Math. 76 (2023), no. 9, 1793--1867.

\bibitem{HanLiu25}
\textsc{J.~Han and Y.~Liu},
On the existence of weighted-cscK metrics,
Adv. Math. 463 (2025), Paper No. 110125, 54 pp.

\bibitem{HanLiu25b}
\textsc{J.~Han and Y.~Liu},
$\mathbb{G}$-uniform weighted K-stability for models on klt varieties,
arXiv:2506.18039v2.

\bibitem{Ha25}
\textsc{Y.~Hashimoto},
Relative uniform K-stability over models implies existence of extremal metrics,
arXiv:2506.02360v1.

\bibitem{HS16}
\textsc{W.~He and S.~Sun},
Frankel conjecture and Sasaki geometry,
Adv. Math. 291 (2016), 912--960.

\bibitem{HN24}
\textsc{T.~Hisamoto and S.~Nakamura},
Continuity method for the Mabuchi soliton on the extremal Fano manifolds,
arXiv:2409.00886, to appear in Ann. Inst. Fourier (Grenoble).

\bibitem{Ino22}
\textsc{E.~Inoue},
Constant {$\mu$}-scalar curvature K\"ahler
metric---formulation and foundational results,
J. Geom. Anal. 32 (2022), no. 5, Paper No. 145, 53 pp.

\bibitem{Jub23}
\textsc{S.~Jubert},
A Yau-Tian-Donaldson correspondence on semisimple principale toric fibrations,
Ann. Inst. Fourier (Grenoble), 73, 2023.

\bibitem{Lah19}
\textsc{A.~Lahdili},
K\"ahler metrics with weighted constant scalar curvature
and weighted K-stability,
Proc. London Math. Soc., 119:1065--114, 2019.

\bibitem{CL21}
\textsc{C.~Li},
Notes on weighted K\"ahler-Ricci solitons
and application to Ricci-flat K\"ahler cone metrics,
arXiv:2107.02088.

\bibitem{CL22}
\textsc{C.~Li},
Geodesic rays and stability in the cscK problem,
Ann. Sci. \'Ec. Norm. Sup\'er. (4) 55 (2022), no. 6, 1529--1574.

\bibitem{Ma86}
\textsc{T.~Mabuchi},
K-energy maps integrating Futaki invariants.
Tohoku Math. J. (2) 38 (1986), no. 4, 575--593.

\bibitem{Ma01}
\textsc{T.~Mabuchi},
K\"{a}hler-Einstein metrics for manifolds with nonvanishing Futaki
character,
Tohoku Math. J, (2) 53 (2001), 171--182.

\bibitem{Ma03}
\textsc{T.~Mabuchi},
Multiplier Hermitian structures on K\"{a}hler manifolds, 
Nagoya Math. J. 170 (2003), 73--115.

\bibitem{Ma21}
\textsc{T.~Mabuchi},
Test configurations, stabilities and canonical K\"ahler metrics-complex geometry
by the energy method,
SpringerBriefs Math,
Springer, Singapore, [2021], x+128 pp.

\bibitem{MSY08}
\textsc{D.~Martelli, J.~Sparks, and S.-T.~Yau},
Sasaki-Einstein manifolds and volume minimization,
Comm. Math. Phys. 280 (2008), 611--673.

\bibitem{Nakg19}
\textsc{Y. Nakagawa},
K\"{a}hler-Ricci soliton on certain toric bundles,
Kyushu J. Math. 73 (2019), no. 2, 379--390.

\bibitem{NN24}
\textsc{Y. Nakagawa and S. Nakamura},
Multiplier Hermitian-Einstein metrics on Fano manifolds
of KSM-type,
Tohoku Math. J. (2) 76 (2024), no. 1, 127--152.

\bibitem{NN25}
\textsc{Y.~Nakagawa and S.~Nakamura},
Modified extremal K\"ahler metrics
and multiplier Hermitian-Einstein metric,
J. Geom. Anal. 35, 104 (2025).

\bibitem{TZ00}
\textsc{G.~Tian and X.~Zhu},
Uniqueness of K\"ahler-Ricci solitons.
Acta Math. 184 (2000), no. 2, 271--305.

\end{thebibliography}



\end{document}


